\documentclass[12pt,fontset=none]{ctexart}
\usepackage[a4paper,margin=1in]{geometry}
\usepackage{enumitem}
\usepackage{hyperref}
\usepackage{fancyhdr}
\usepackage{xcolor}
\setlength{\headheight}{15pt}
\setCJKmainfont{STSong}
\setCJKsansfont{STSong}
\setCJKmonofont{STSong}
\pagestyle{fancy}
\fancyhf{}
\lhead{Cell Biology Mock Exam}
\rhead{Based on notebook + PDFs}
\cfoot{\thepage}
\setlist[enumerate]{leftmargin=*}

\title{Cell Biology Mock Exam and Answer Key}
\author{}
\date{2026-04-27}

\begin{document}
\maketitle

\section*{Exam}
\textbf{Instructions}
\begin{itemize}
\item This mock exam is designed to roughly match the style and difficulty of \texttt{2026CB\_homework\_I.pdf}.
\item It aims to cover about 70--80\% of the major knowledge points in the current notebook and lecture PDFs.
\item Total score: 102.
\end{itemize}

\subsection*{I. Multiple choice questions}
\textbf{Questions 1--62, 1 point each. Select the best answer unless otherwise instructed.}

\begin{enumerate}
\item Which of the following is the best example of a prokaryotic homolog of tubulin?
\begin{enumerate}[label=\alph*.]
\item MreB
\item ParM
\item FtsZ
\item Crescentin
\end{enumerate}

\item Which cytoskeletal system is most directly responsible for the contractile ring during cytokinesis?
\begin{enumerate}[label=\alph*.]
\item Actin filaments with myosin II
\item Intermediate filaments
\item Microtubules with dynein
\item Centrosomal lamins
\end{enumerate}

\item Which cytoskeletal element is most important for long-distance intracellular transport of membrane-enclosed organelles?
\begin{enumerate}[label=\alph*.]
\item Actin cortex
\item Microtubules
\item Intermediate filaments
\item Desmosomes
\end{enumerate}

\item Stress fibers are best described as:
\begin{enumerate}[label=\alph*.]
\item Branched actin networks under the plasma membrane
\item Contractile actin bundles in non-muscle cells
\item Intermediate filament cables in epithelia
\item Microtubule bundles in migrating cells
\end{enumerate}

\item Which protein complex primarily nucleates branched actin networks?
\begin{enumerate}[label=\alph*.]
\item Formin
\item Arp2/3 complex
\item Dynactin
\item Katanin
\end{enumerate}

\item Which protein is most directly associated with elongation of long, unbranched actin filaments?
\begin{enumerate}[label=\alph*.]
\item Formin
\item Arp2/3
\item CapZ
\item Cofilin
\end{enumerate}

\item CapZ primarily acts by:
\begin{enumerate}[label=\alph*.]
\item Severing actin filaments
\item Binding ATP-actin monomers
\item Capping the plus end of actin filaments
\item Linking dynein to membranes
\end{enumerate}

\item Cofilin promotes actin filament turnover mainly by:
\begin{enumerate}[label=\alph*.]
\item Capping the minus end
\item Twisting and destabilizing the filament
\item Converting ADP to ATP on actin
\item Anchoring actin to Z discs
\end{enumerate}

\item The motor activity of myosin II is most closely coupled to:
\begin{enumerate}[label=\alph*.]
\item GTP hydrolysis
\item ATP hydrolysis and Pi release
\item ADP synthesis
\item Ca$^{2+}$ export
\end{enumerate}

\item In skeletal muscle, the plus ends of thin filaments are anchored at the:
\begin{enumerate}[label=\alph*.]
\item M line
\item A band
\item Z disc
\item H zone
\end{enumerate}

\item Which protein acts as a molecular spring in the sarcomere?
\begin{enumerate}[label=\alph*.]
\item Tropomodulin
\item Titin
\item CapZ
\item Dynein
\end{enumerate}

\item Smooth muscle contraction is directly activated when Ca$^{2+}$:
\begin{enumerate}[label=\alph*.]
\item Binds troponin C on actin
\item Binds calmodulin and activates MLCK
\item Causes CapZ dissociation
\item Stabilizes tropomyosin on microtubules
\end{enumerate}

\item Dynamic instability of microtubules depends primarily on:
\begin{enumerate}[label=\alph*.]
\item ATP hydrolysis on $\alpha$-tubulin
\item GTP hydrolysis on $\beta$-tubulin
\item Dynein detachment
\item Actin treadmilling
\end{enumerate}

\item Loss of the GTP cap from a microtubule plus end usually triggers:
\begin{enumerate}[label=\alph*.]
\item Rescue
\item Treadmilling
\item Catastrophe
\item Annealing
\end{enumerate}

\item The major function of $\gamma$-tubulin is to:
\begin{enumerate}[label=\alph*.]
\item Cap actin filaments
\item Nucleate microtubules
\item Stabilize intermediate filaments
\item Carry cargo to the Golgi
\end{enumerate}

\item Which protein severs microtubules?
\begin{enumerate}[label=\alph*.]
\item Katanin
\item Tropomodulin
\item Lis1
\item Ankyrin
\end{enumerate}

\item Most kinesins move toward the:
\begin{enumerate}[label=\alph*.]
\item Minus end of actin filaments
\item Plus end of microtubules
\item Minus end of microtubules
\item Centrosomal matrix
\end{enumerate}

\item Dynein attachment to membrane-enclosed organelles is commonly mediated by:
\begin{enumerate}[label=\alph*.]
\item CENP-A
\item Dynactin
\item H1
\item SRP
\end{enumerate}

\item In motile cilia and flagella, dynein-driven sliding between adjacent doublets is converted to bending by:
\begin{enumerate}[label=\alph*.]
\item DNA looping
\item Flexible linker proteins
\item ATP synthase
\item Capping proteins
\end{enumerate}

\item Plectin is best described as a protein that:
\begin{enumerate}[label=\alph*.]
\item Generates ATP in mitochondria
\item Cross-links intermediate filaments to other cytoskeletal systems
\item Caps the plus end of actin filaments
\item Forms the nuclear pore basket
\end{enumerate}

\item Members of the nuclear lamina structural network belong to which family?
\begin{enumerate}[label=\alph*.]
\item Type V intermediate filaments
\item Actin-related proteins
\item Kinesin family proteins
\item Histone variants
\end{enumerate}

\item Which one is an A-type lamin?
\begin{enumerate}[label=\alph*.]
\item Lamin B1
\item Lamin B2
\item Lamin C
\item CENP-A
\end{enumerate}

\item Karyotyping is most useful for detecting:
\begin{enumerate}[label=\alph*.]
\item Point mutations in coding regions
\item Single amino acid substitutions
\item Aneuploidy and large chromosomal rearrangements
\item mRNA splicing defects
\end{enumerate}

\item Which statement about nucleosome positioning is most accurate?
\begin{enumerate}[label=\alph*.]
\item DNA sequence has no influence on nucleosome formation
\item Only histone H1 determines nucleosome location
\item DNA bendability contributes to preferred nucleosome positions
\item Nucleosomes form only on repetitive DNA
\end{enumerate}

\item Linker histone H1 primarily:
\begin{enumerate}[label=\alph*.]
\item Replaces histone H3 in centromeres
\item Changes DNA exit path and promotes chromatin compaction
\item Catalyzes DNA methylation
\item Binds directly to RNA polymerase II CTD
\end{enumerate}

\item ATP-dependent chromatin remodeling complexes can:
\begin{enumerate}[label=\alph*.]
\item Translate mRNA into protein
\item Slide, remove, or exchange nucleosomes
\item Hydrolyze histones into peptides
\item Ligate Okazaki fragments
\end{enumerate}

\item The histone variant that helps define many centromeres is:
\begin{enumerate}[label=\alph*.]
\item H1
\item HP1
\item CENP-A
\item Xist
\end{enumerate}

\item Self-propagation of heterochromatin is best explained by:
\begin{enumerate}[label=\alph*.]
\item Random diffusion of histones
\item Reader-writer positive feedback
\item Ribosome assembly
\item GPI anchoring
\end{enumerate}

\item Which method gives a population-averaged contact map of chromosome conformation?
\begin{enumerate}[label=\alph*.]
\item FRAP
\item 3C/Hi-C
\item FRET
\item Patch clamp
\end{enumerate}

\item Chromosome tracing is especially valuable because it:
\begin{enumerate}[label=\alph*.]
\item Measures ATP turnover in chromatin
\item Directly visualizes chromosome organization in single cells
\item Replaces all sequencing-based methods
\item Detects only coding sequences
\end{enumerate}

\item In polytene chromosomes, the chromocenter forms from:
\begin{enumerate}[label=\alph*.]
\item Mitochondrial DNA clusters
\item Aggregated heterochromatin near centromeres
\item Fused nucleoli
\item Centrosomal satellites
\end{enumerate}

\item Gene duplication is evolutionarily important because it:
\begin{enumerate}[label=\alph*.]
\item Prevents any future mutation
\item Provides raw material for new gene functions
\item Eliminates homologous genes
\item Always reduces genome size
\end{enumerate}

\item Which RNA polymerase mainly transcribes mRNAs?
\begin{enumerate}[label=\alph*.]
\item Pol I
\item Pol II
\item Pol III
\item Pol IV
\end{enumerate}

\item Which RNA polymerase mainly transcribes tRNAs and 5S rRNA?
\begin{enumerate}[label=\alph*.]
\item Pol I
\item Pol II
\item Pol III
\item Telomerase
\end{enumerate}

\item DNA superhelical tension generated during transcription is relieved by:
\begin{enumerate}[label=\alph*.]
\item Topoisomerases
\item Restriction enzymes
\item SRP
\item Cohesin
\end{enumerate}

\item The earliest processing event on a eukaryotic pre-mRNA is:
\begin{enumerate}[label=\alph*.]
\item Polyadenylation
\item 5' capping
\item Spliceosome disassembly
\item tRNA charging
\end{enumerate}

\item Formation of the 3' end of most eukaryotic mRNAs requires:
\begin{enumerate}[label=\alph*.]
\item RNA editing only
\item Direct transcription termination with no processing
\item Cleavage followed by poly(A) addition
\item Reverse transcription
\end{enumerate}

\item DNA looping in gene regulation is important because it:
\begin{enumerate}[label=\alph*.]
\item Prevents all enhancer action
\item Allows distant regulators to contact promoter-bound machinery
\item Converts DNA into RNA
\item Forces nucleosomes off the chromosome permanently
\end{enumerate}

\item Transcriptional synergy refers to:
\begin{enumerate}[label=\alph*.]
\item Equal contribution of all genes
\item A less-than-additive effect of two activators
\item A greater-than-additive effect of multiple activators
\item Inhibition of one activator by another
\end{enumerate}

\item Insulator DNA sequences primarily function to:
\begin{enumerate}[label=\alph*.]
\item Recruit ribosomes
\item Block inappropriate long-range gene regulation
\item Replicate telomeres
\item Polymerize actin
\end{enumerate}

\item A major molecular basis for cell memory is:
\begin{enumerate}[label=\alph*.]
\item Positive feedback and heritable chromatin states
\item Loss of all transcription factors after mitosis
\item Removal of all DNA methylation
\item Complete absence of signaling pathways
\end{enumerate}

\item Genomic imprinting is best defined as:
\begin{enumerate}[label=\alph*.]
\item Expression of both alleles equally
\item Parent-of-origin-dependent monoallelic expression
\item Random loss of one chromosome
\item Universal repression of maternal genes
\end{enumerate}

\item X-chromosome inactivation in mammals is most closely associated with:
\begin{enumerate}[label=\alph*.]
\item CapZ
\item Xist RNA
\item Tropomodulin
\item SRP54
\end{enumerate}

\item Alternative cleavage and polyadenylation (APA) can change:
\begin{enumerate}[label=\alph*.]
\item Only the DNA replication origin
\item The fate and coding output of the transcript
\item Ribosome structure
\item The polarity of actin filaments
\end{enumerate}

\item The most abundant phospholipid in many eukaryotic membranes is:
\begin{enumerate}[label=\alph*.]
\item Phosphatidylcholine
\item Cardiolipin
\item Phosphatidylinositol trisphosphate
\item Cholesteryl ester
\end{enumerate}

\item Cholesterol helps membranes by:
\begin{enumerate}[label=\alph*.]
\item Hydrolyzing phospholipids
\item Filling gaps created by unsaturated hydrocarbon tails
\item Removing all membrane asymmetry
\item Serving as a motor protein
\end{enumerate}

\item Which membrane motion is rare without enzymes?
\begin{enumerate}[label=\alph*.]
\item Lateral diffusion
\item Rotation
\item Flexion
\item Flip-flop
\end{enumerate}

\item Which class of membrane proteins usually requires detergent for solubilization?
\begin{enumerate}[label=\alph*.]
\item Peripheral membrane proteins
\item Integral membrane proteins
\item Histones
\item Ribosomal proteins
\end{enumerate}

\item Hydropathy plots are most useful for predicting:
\begin{enumerate}[label=\alph*.]
\item DNA methylation islands
\item Potential membrane-spanning hydrophobic segments
\item RNA splicing enhancers
\item Chromosome territories
\end{enumerate}

\item Glycosylation of membrane proteins generally faces the:
\begin{enumerate}[label=\alph*.]
\item Cytosolic side
\item Non-cytosolic/extracellular side
\item Centrosomal side
\item Nucleoplasmic side
\end{enumerate}

\item Which of the following is NOT one of the three major membrane-bending mechanisms emphasized in class?
\begin{enumerate}[label=\alph*.]
\item Hydrophobic insertion
\item Scaffolding
\item Protein crowding
\item DNA ligation
\end{enumerate}

\item Transporters differ from channels mainly because transporters:
\begin{enumerate}[label=\alph*.]
\item Never bind cargo
\item Undergo alternating-access conformational changes
\item Are always faster than channels
\item Can move only water
\end{enumerate}

\item The three classic states of a transporter are:
\begin{enumerate}[label=\alph*.]
\item Bound, unbound, degraded
\item Open, closed, cleaved
\item Outward-open, occluded, inward-open
\item Resting, spiking, refractory
\end{enumerate}

\item Which of the following is NOT a major class of ATP-driven pump?
\begin{enumerate}[label=\alph*.]
\item P-type
\item ABC
\item F-type/V-type
\item Arp-type
\end{enumerate}

\item Bacterial cells protect themselves against extreme osmotic shock using:
\begin{enumerate}[label=\alph*.]
\item Mechanosensitive channels
\item Histone H1
\item Nuclear pores
\item Tight junctions
\end{enumerate}

\item The nicotinic acetylcholine receptor at the neuromuscular junction is:
\begin{enumerate}[label=\alph*.]
\item A GPCR
\item A voltage-gated K$^+$ channel
\item A ligand-gated cation channel
\item A proton pump
\end{enumerate}

\item Long-term potentiation in the hippocampus depends critically on:
\begin{enumerate}[label=\alph*.]
\item Ca$^{2+}$ entry through NMDA receptors
\item ATP export through VDAC
\item Histone H1 phosphorylation
\item Dynein-driven vesicle bending
\end{enumerate}

\item Which transport mechanism moves folded proteins between nucleus and cytosol?
\begin{enumerate}[label=\alph*.]
\item Vesicular transport
\item Gated transport
\item Transmembrane translocation through TOM
\item Endocytosis
\end{enumerate}

\item The outer nuclear membrane is continuous with the:
\begin{enumerate}[label=\alph*.]
\item Lysosomal membrane
\item ER membrane
\item Plasma membrane
\item Mitochondrial inner membrane
\end{enumerate}

\item Most proteins imported into mitochondria enter first through:
\begin{enumerate}[label=\alph*.]
\item Sec61
\item TOM complex
\item NPC
\item COPII
\end{enumerate}

\item A classic peroxisomal targeting signal is:
\begin{enumerate}[label=\alph*.]
\item KDEL at the N-terminus
\item SKL at the C-terminus
\item NLS rich in lysine
\item A transmembrane $\beta$-barrel
\end{enumerate}

\item Which statement best describes the role of SRP?
\begin{enumerate}[label=\alph*.]
\item It splices introns from mRNA
\item It recognizes ER signal peptides and pauses translation transiently
\item It degrades misfolded proteins in lysosomes
\item It methylates histones at centromeres
\end{enumerate}
\end{enumerate}

\subsection*{II. Short-answer questions}
\textbf{Questions 1--8, 5 points each.}

\begin{enumerate}
\item A dividing cell is treated with nocodazole, while another is treated with Taxol. Both populations stop proliferating. Explain why these two drugs, despite having opposite effects on microtubule polymer mass, can both block mitosis.

\item Compare branched actin networks and actin bundles in terms of nucleating proteins, architecture, and typical cellular functions.

\item Explain why membrane asymmetry is functionally important, and why the orientation of membrane proteins and lipids is preserved during vesicular transport.

\item A mutation disrupts the reader protein in a heterochromatin spreading system. Predict the effect on heterochromatin propagation and explain the logic using the reader-writer model.

\item Explain how Ran GTPase imposes directionality on nuclear import and export.

\item A membrane protein is correctly synthesized but persistently misfolds in the ER. Describe the pathway that removes it and explain why this pathway is important for cell physiology.

\item Compare 3C/Hi-C and chromosome tracing as methods for studying 3D genome organization. What does each method measure best, and why are they complementary?

\item Explain how alternative cleavage and polyadenylation could allow the same gene to produce a membrane-bound form in one cell type and a secreted form in another.
\end{enumerate}

\newpage
\section*{Answer Key}

\subsection*{I. Multiple choice answers}
\begin{enumerate}
\item c
\item a
\item b
\item b
\item b
\item a
\item c
\item b
\item b
\item c
\item b
\item b
\item b
\item c
\item b
\item a
\item b
\item b
\item b
\item b
\item a
\item c
\item c
\item c
\item b
\item b
\item c
\item b
\item b
\item b
\item b
\item b
\item c
\item c
\item a
\item b
\item c
\item b
\item c
\item b
\item a
\item b
\item b
\item b
\item a
\item b
\item d
\item b
\item b
\item d
\item d
\item b
\item c
\item a
\item c
\item a
\item b
\item b
\item b
\item b
\item b
\item b
\end{enumerate}

\subsection*{II. Short-answer reference answers}
\begin{enumerate}
\item \textbf{Nocodazole} promotes microtubule depolymerization, whereas \textbf{Taxol} stabilizes microtubules and suppresses their dynamic instability. Mitosis requires highly dynamic spindle microtubules to search for kinetochores, capture chromosomes, align them, and then reorganize during segregation. Therefore, both too little microtubule polymer and overly stabilized, non-dynamic microtubules disrupt spindle function and block mitosis.

\item Branched actin networks are nucleated mainly by the \textbf{Arp2/3 complex}, which generates daughter filaments from the sides of preexisting filaments, producing dense dendritic arrays. These networks are typical of lamellipodia and the cell cortex, where they help drive membrane protrusion and rapid shape change. In contrast, actin bundles are commonly built from long, relatively straight filaments promoted by \textbf{formin}; they are then organized by cross-linking proteins into parallel or contractile bundles. Bundles function in structures such as microvilli, filopodia, stress fibers, and the contractile ring.

\item Membrane asymmetry is important because the two leaflets and the two faces of membrane proteins perform different functions. For example, glycolipids and many glycosylated domains face the non-cytosolic/extracellular side, whereas many signaling and cytoskeletal interactions occur on the cytosolic side. During vesicular transport, membrane topology is preserved: the lumen of the ER/Golgi corresponds topologically to the extracellular space, so proteins and lipids do not flip their orientation when they move in vesicles.

\item In the reader-writer model, the \textbf{reader} recognizes an existing heterochromatin mark and recruits the \textbf{writer}, which deposits the same mark on adjacent nucleosomes. If the reader is lost, the system can no longer efficiently recognize the existing modification and spread it outward. As a result, heterochromatin propagation is impaired or stops, and the self-propagating maintenance of the silent state becomes unstable.

\item Ran directionality depends on compartment-specific regulation of its nucleotide state. In the nucleus, \textbf{Ran-GEF} generates high \textbf{Ran-GTP}; in the cytosol, \textbf{Ran-GAP} promotes Ran-GTP hydrolysis to Ran-GDP. During import, import receptors bind cargo in the cytosol, pass through the nuclear pore, and release cargo when they encounter Ran-GTP in the nucleus. During export, export receptors bind cargo together with Ran-GTP in the nucleus, carry the complex out, and then release cargo after GTP hydrolysis in the cytosol. This asymmetric Ran system imposes directionality on both import and export.

\item Persistently misfolded ER proteins are removed by \textbf{ER-associated degradation (ERAD)}. The protein is recognized by the ER quality-control system, retrotranslocated or dislocated back across the ER membrane, ubiquitinated, extracted by ATP-dependent factors, and degraded by the proteasome in the cytosol. ERAD is important because it prevents toxic accumulation of misfolded proteins, maintains ER homeostasis, and helps limit ER stress.

\item \textbf{3C/Hi-C} measures contact frequencies between chromosomal loci across populations of cells, making it especially powerful for identifying average features such as compartments, TAD-like organization, and long-range interaction tendencies. \textbf{Chromosome tracing} uses imaging to visualize chromosome folding directly in single cells, so it is best for revealing cell-to-cell variability and actual spatial configurations. They are complementary because sequencing-based contact maps provide global statistical organization, whereas imaging reveals single-cell structural heterogeneity.

\item Alternative cleavage and polyadenylation can place the 3' end of an mRNA at different positions. If a proximal poly(A) site is used upstream of sequences encoding a transmembrane domain, the resulting mRNA may exclude that membrane-anchor coding region and produce a \textbf{secreted form}. If a distal poly(A) site is used, the transcript may retain the transmembrane segment and produce a \textbf{membrane-bound form}. Thus, cell-type-specific APA can alter both transcript structure and protein destination.
\end{enumerate}

\end{document}
